%!TEX encoding = UTF-8 Unicode
%!TEX root = ../lect-w07.tex

\ifkompendium\else
\Subsection{Grumligtlådan}

\begin{Slide}{Grumligtlådan: topplista med ämnen}
\begin{multicols}{2}
\begin{verbatim}
  om kursen       : 14
  klass           : 9
  this            : 7
  filstruktur     : 5
  pluggteknik     : 4
  case-klass      : 4
  registrera      : 3
  filtrera        : 3
  copy            : 3
  tupel           : 2
  konstruktor     : 2
  iterera         : 2
  fabriksmetod    : 2
  synlighet       : 1
  skuggning       : 1
  sidoeffekt      : 1
  Seq[T]          : 1
  sekvenser       : 1
  sats/uttryck    : 1
  lambda          : 1
  kompanjonsobj   : 1
  högre ordn funk : 1
  funktion        : 1
  curry-funktion  : 1
\end{verbatim}
\end{multicols}
\end{Slide}



\begin{Slide}{Grumligtlådan: lajvkodning}
Lajvkodning enligt begrepp i lådan utifrån denna klass:
\begin{Code}[basicstyle=\ttfamily\SlideFontSize{12}{14}]
class Frog(var n: Int) { def hop = n += 1 }
\end{Code}
\includegraphics[width=0.6\textwidth]{../img/frog}
\end{Slide}

\fi
